%---------------------------------------------------------------------------%
%->> Main content
%---------------------------------------------------------------------------%
\clearpage
\section{选题的背景及意义}

\subsection{选题背景}

单细胞RNA测序（single-cell RNA sequencing, scRNA-seq）技术的快速发展，使得研究人员能够在单个细胞层面解析基因表达的异质性，为理解复杂生物系统的细胞组成和功能状态提供了非常高的分辨率\cite{kolodziejczyk2015technology}。近年来，随着高通量测序成本的持续下降和实验方法的不断成熟，单细胞转录组学已经成为现代生物学研究的核心技术之一，在发育生物学、免疫学、肿瘤学以及神经科学等多个领域发挥着关键作用\cite{tanay2017scaling}。然而，传统的单细胞测序技术本质上是一种观测性（observational）手段，仅能刻画细胞在特定时间点的静态表达状态，难以揭示基因调控的因果关系和细胞状态转换的动态机制。

为了深入理解基因功能和细胞命运决定的分子机制，扰动实验（perturbation experiments）应运而生。通过人为干预特定基因的表达水平，例如利用CRISPR-Cas9基因编辑系统进行基因敲除（knockout）、基因激活（activation）或干扰（interference），研究人员能够直接观察基因扰动对细胞转录组的影响，从而建立基因间的因果调控关系\cite{dixit2016perturb}。这种将扰动与单细胞测序相结合的技术，即Perturb-seq，能够在单细胞分辨率下同时读取扰动信息和转录组响应，为系统级的基因功能注释和调控网络推断提供了强大的实验范式\cite{adamson2016multiplexed}。

然而，大规模扰动实验面临着诸多挑战。首先，实验成本高昂：一次完整的Perturb-seq实验通常需要对数千至数万个细胞进行测序，在全基因组尺度上系统性地筛选所有基因的扰动效应仍然耗资巨大\cite{replogle2022mapping}。其次，实验周期较长：从细胞培养、病毒感染、药物处理到文库构建和测序分析，整个流程通常需要数周时间。此外，某些扰动条件在实验中难以实现或存在技术限制，例如特定的多基因联合扰动、剂量依赖的连续扰动以及时间序列的动态扰动响应等\cite{norman2019exploring}。这些因素共同制约了扰动数据的规模化获取，使得当前可用的扰动数据集在扰动类型和细胞数量上都存在显著的覆盖缺口。

在此背景下，计算方法应运而生，旨在利用有限的观测数据和扰动数据，预测未见过的扰动条件下的细胞转录组响应\cite{rood2024toward, lotfollahi2023contextualization}。这种计算预测范式具有重要的应用价值：一方面，它可以大幅降低实验成本，通过预测筛选优先关注最具生物学意义的扰动条件；另一方面，它能够扩展扰动空间的覆盖范围，预测实验难以实现的复杂扰动组合或条件。近年来，基于深度学习的扰动预测方法取得了显著进展，各种神经网络架构被应用于建模基因表达的高维分布和扰动响应的非线性映射关系\cite{lotfollahi2019scgen, lotfollahi2022predicting}。

在深度生成模型领域，扩散模型（diffusion models）和流匹配（flow matching）等方法在图像生成、语音合成等任务中取得了突破性进展\cite{ho2020denoising, song2020score}。其中，条件流匹配（Conditional Flow Matching, CFM）作为一种新兴的生成模型框架，通过学习从噪声分布到目标数据分布的连续传输映射，展现出在高维复杂数据建模中的独特优势\cite{lipman2022flow, liu2022flow}。与传统的扩散模型相比，条件流匹配具有训练稳定性高、采样效率好、理论基础清晰等优点，为解决单细胞扰动预测问题提供了新的技术路径\cite{huguet2024cellflow}。

与此同时，因果表示学习（causal representation learning）领域的进展为扰动预测提供了新的理论视角\cite{scholkopf2021toward}。单细胞扰动数据本质上蕴含了丰富的因果信息：扰动作为一种干预，使得观测数据包含了因果结构的线索。Lorch等人提出的潜在因果扩散模型（Latent Causal Diffusion, LCD）将基因表达建模为稳态扩散过程，通过学习基因调控的动态系统来预测扰动响应，并发展了CLIPR方法从学习到的动力学中提取可解释的因果效应\cite{lorch2025latent}。Jiang等人提出的PerturbedVAE框架从扰动抑制假说出发，揭示了扰动特定信息在基因表达数据中的稀疏性本质，通过显式分离扰动不变信息和扰动响应信息，实现了具有理论可识别性保证的因果表示学习\cite{jiang2025what}。这些工作为将因果推理与深度生成模型相融合提供了重要的方法论基础。

\subsection{选题意义}

单细胞扰动预测本质上是一个条件生成与因果推断交叉的问题：模型既需要从有限实验数据中生成符合真实生物分布的扰动后转录组，也需要理解扰动如何改变细胞状态并识别少数真正响应的关键基因程序。然而，当前方法面临一个被忽视的技术困境——单细胞转录组维度高（通常数千至数万个基因），而真实扰动响应集中在一小部分基因上，大量无响应基因在损失函数加和过程中淹没了响应基因的监督信号，导致模型倾向于预测全局近零变化来换取较低的整体MSE，差异表达基因的恢复因此变得不稳定。此外，现有方法普遍将扰动基因编码为独立嵌入，未能充分利用基因功能关联（GO）、蛋白互作网络（STRING）以及转录调控关系（如TRRUST）等多层次先验信息，模型在扰动基因间的协同效应建模和下游靶基因识别方面受到限制。

针对上述问题，本课题从两个互补维度展开研究。其一，提出基于先验约束的稀疏响应扰动预测方法，将GO功能关联、STRING蛋白互作和TRRUST转录调控关系三层先验信息融入条件流匹配框架——通过图注意力网络编码扰动基因间的功能与调控上下文，借助条件到基因的交叉注意力精确检索受影响的基因程序，并利用扰动条件感知的gene mask和SNR加权损失在基因粒度上施加稀疏响应约束，引导模型聚焦少数高可信响应基因。其二，探索基于因果解耦的扰动表示学习，将基因表达显式分离为扰动不变的背景变量和扰动响应的因果变量，使模型区分”细胞本身的背景差异”和”扰动真正诱导的变化”，提升跨条件预测的泛化性和可解释性。

应用价值方面，精准的扰动预测模型在生物医药领域具有广泛前景。在药物研发方面，模型可用于虚拟筛选候选药物或基因靶点的转录组效应，减少高成本实验筛选的规模\cite{tejada2025causal}；在疾病机制研究方面，通过模拟疾病相关基因的功能缺失或获得，有助于揭示分子病理机制和潜在治疗靶点\cite{uhler2024building}；在合成生物学领域，扰动预测可辅助设计基因回路和优化工程化细胞系统。总体而言，本课题通过”先验约束驱动稀疏建模”与”因果解耦表征学习”两个研究点，既提升虚拟扰动预测对关键响应基因和扰动方向的刻画能力，也增强模型对扰动机制的解释能力。

\clearpage
\section{国内外本学科领域的发展现状与趋势}

\subsection{单细胞扰动发展现状}

单细胞扰动预测旨在利用计算模型预测基因或化学扰动对细胞转录组的影响，是近年来计算生物学领域的研究热点。不同方法在建模范式、先验知识利用和生成质量上各有侧重，以下从数据驱动的生成建模和知识引导与因果驱动两个维度进行梳理。

\subsubsection{数据驱动的生成式扰动建模}

早期方法以线性建模和自编码器为主。scGen\cite{lotfollahi2019scgen}将细胞投影至VAE潜在空间，在潜在空间中用线性扰动向量建模扰动效应。CPA（Compositional Perturbation Autoencoder）\cite{lotfollahi2023contextualization}将其扩展为组合式框架，支持多药物、多剂量、多细胞类型建模，能够从单基因扰动数据外推组合扰动效应。SAMS-VAE\cite{bereket2023sams}和sVAE+\cite{lopez2022svae}在潜在空间中引入稀疏结构或机制偏移来区分不同扰动类型的影响。这些方法计算高效，但线性或浅层非线性假设限制了复杂扰动响应的刻画。

在分布映射方面，CellOT\cite{bunne2023learning}基于最优传输（Optimal Transport）理论，将扰动预测形式化为求解Kantorovich问题或其熵正则化形式，学习从对照分布到扰动分布的传输映射。Waddington-OT\cite{schiebinger2019optimal}利用最优传输重建细胞发育动态轨迹。扩散模型方面，PerturbDiff\cite{roman2024perturbdiff}通过随机微分方程（SDE）建模扰动响应，SCDFM\cite{wu2025scdfm}进一步扩展至多批次整合和伪时间分析。CellFlow\cite{huguet2024cellflow}首次将条件流匹配（Conditional Flow Matching, CFM）\cite{lipman2022flow}应用于扰动预测，通过学习条件向量场构建从对照到扰动状态的确定性传输映射，在多个基准上超越CPA等方法且推理速度提升20倍以上。与扩散模型相比，CFM直接回归向量场而无需求解得分匹配，采样路径确定性，在高维生成任务中展现出训练稳定性和采样效率优势\cite{liu2022flow}。

Transformer架构和预训练范式也推动了该领域发展。scGPT\cite{cui2024scgpt}和scBERT\cite{yang2022scbert}在大规模单细胞数据上预训练，学习细胞转录组的通用表示后在下游扰动任务上微调。Geneformer\cite{theodoris2023transferable}将基因表达值转化为排序序列，利用Transformer学习基因间上下文依赖。SCLang\cite{zeng2025sclang}探索了转录组、细胞生物学和自然语言的多模态对齐。值得注意的是，近期研究指出基础模型的通用预训练表示在扰动标签的线性可解码性上甚至弱于PCA基线\cite{jiang2025what}，提示扰动相关信息在通用预训练中未必能被充分保留。

\subsubsection{知识引导与因果驱动的扰动建模}

生物学先验知识——包括基因功能注释（Gene Ontology）、蛋白-蛋白互作网络（STRING、BioGRID）以及转录调控关系（TRRUST\cite{han2018trrust}、DoRothEA\cite{garcia2019dorothea}、RegNetwork）——为高维稀疏转录组数据提供了有价值的归纳偏置。GEARS\cite{roohani2024gears}利用基因共表达和GO信息构建基因关联图，通过图卷积网络学习扰动响应。TxPert\cite{wu2026txpert}进一步整合STRING、GO和扰动筛选衍生图等多种知识图谱，利用异构图编码器学习扰动基因或基因产物间的多维度关系，说明了多层次知识图谱在扰动条件编码中的互补价值。在更广泛的计算生物学领域，图神经网络已被广泛应用于基因功能预测\cite{zhang2022graph}、药物-靶标相互作用预测和蛋白质结构分析等任务，但如何最有效地将异质先验图嵌入条件生成模型的各个环节——从条件编码、特征融合到训练目标设计——仍是活跃的研究前沿。

从基因调控网络（GRN）的角度，基因扰动实验本质上是对GRN的定向干预。TRRUST v2\cite{han2018trrust}收录了约800个人类转录因子及其约8400条调控关系，全部来源于文献人工整理，置信度较高。DoRothEA\cite{garcia2019dorothea}按证据强度分级提供TF regulon。然而，目前绝大多数扰动预测方法并未显式利用GRN中的TF→靶基因调控方向信息，而是将其与无向的功能关联或互作网络等同处理，错失了调控方向这一关键先验。

因果表示学习为该领域提供了从"拟合数据分布"走向"理解扰动机制"的方法论路径\cite{scholkopf2021toward}。Discrepancy-VAE\cite{zhang2023identifiability}通过利用不同扰动条件下的分布差异，实现了潜在因果变量的可识别学习。CausalCell\cite{deng2024causal}和CausCell\cite{gao2025causal}分别通过因果图和结构因果模型（SCM）引入因果结构先验。LCD\cite{lorch2025latent}将基因表达建模为由SDE控制的稳态扩散过程，扰动被视为对底层基因调控机制的修改，通过两阶段推断从稳态快照中学习漂移函数，在预测未见扰动组合的分布偏移方面显著优于纯数据驱动方法，CLIPR方法进一步从学习到的动力学中提取基因间直接因果效应。PerturbedVAE\cite{jiang2025what}从扰动抑制假说出发，揭示了基因表达以扰动不变信息为主导的挑战，通过潜在因果生成模型将扰动不变变量$z_\iota$与扰动响应变量$z_\nu$显式分离，理论分析给出了在部分干预数据下可靠恢复因果变量的条件。

\subsection{单细胞扰动研究趋势}

从方法演进角度看，单细胞扰动预测正经历三个层面的深化。在生成能力层面，从早期的VAE线性模型到扩散模型和条件流匹配，生成质量和采样效率持续提升。条件流匹配因无需得分匹配、采样确定性、支持灵活的条件路径设计等特性，在高维转录组生成中展现出显著潜力。在知识融合层面，先验生物学知识正从辅助特征升级为核心建模组件——从GEARS利用GO构建简单关联图，到TxPert融合多种知识图谱的异质图编码，再到将TRRUST等有向调控网络作为基因级注意力偏置引导模型聚焦已知调控靶点，多层次先验的系统嵌入正在改变纯数据驱动方法在高维稀疏数据上的困境。在机制理解层面，因果表示学习的引入使模型不仅能够预测扰动响应，还能区分背景变异与真实的因果效应，其中LCD从稳态快照推断动态调控机制、PerturbedVAE通过扰动抑制假说验证因果解耦的必要性，代表了该方向的方法论深化。

在任务维度上，领域正从单基因扰动向组合扰动、从稳态快照向动态过程、从单一转录组向多组学整合扩展。LCD将基因调控形式化为随机微分方程控制的动态系统\cite{lorch2025latent}，PerturbedVAE在单基因到双基因扰动外推上展现了因果表示学习的优势\cite{jiang2025what}。总体而言，多层次先验知识引导、因果解耦表征与高效生成模型的深度融合，构成了当前单细胞扰动预测领域的主要发展方向，也是本课题的核心研究思路。

\clearpage
\section{课题主要研究内容、预期目标}

针对当前单细胞扰动预测方法在生成质量、基因级扰动响应刻画和因果可解释性方面的不足，本课题拟从两个互补的研究维度展开：

\subsection{基于先验增强的稀疏响应扰动预测}

本研究的第一个维度聚焦于将多层次生物先验与稀疏响应约束系统融入条件流匹配框架，重点解决高维转录组扰动预测中无响应基因梯度淹没响应基因监督信号、差异表达基因恢复不稳定的问题。

\textbf{问题背景。}单细胞扰动预测的核心任务是：给定一个基因扰动条件（如”敲除基因A”），预测细胞转录组的全局变化。当前基于条件流匹配（Conditional Flow Matching, CFM）的深度生成模型在此任务上取得了显著进展，但在基因粒度的预测精度上面临一个根本困境——单细胞转录组包含上万个基因，而每个扰动在生物学上仅真正影响少数关键基因或通路。现有方法以均方误差（MSE）作为损失函数，对所有基因等权加和，导致大量无响应基因的梯度在优化过程中占据主导，淹没了少数真正响应基因的监督信号。其直接后果是模型倾向于预测”全局近零变化”来换取较低的整体MSE，差异表达基因（DEG）的恢复极不稳定。换句话说，模型的整体MSE可以很低，但真正重要的那几十个响应基因的预测方向可能是错的——这是一个典型的”评价指标与生物目标错配”问题。

\textbf{现有方法的局限。}现有方法（如CPA、CellFlow、scDFM等）在处理扰动条件时普遍采用one-hot向量或静态嵌入来表示扰动基因，与表达向量简单拼接后输入网络。这种设计在基因粒度上存在两个盲区：其一，扰动基因被当作匿名标签处理，丢失了基因间的功能关联（Gene Ontology）、蛋白物理互作（STRING）以及转录调控方向（TRRUST）等多层次先验信息——模型不知道”哪些基因在功能上相似、在通路中相邻、在调控上联通”，因此无法利用这些先验来推断一个扰动可能影响哪些下游基因；其二，损失函数在基因维度上的无差别加和，使得模型在训练时没有内在机制将有限的模型容量优先分配给少数关键响应基因。GEARS等方法开始利用GO和基因共表达图来增强扰动表示，TxPert进一步融合了多种知识图谱，但这些方法尚未将多层次异质先验与稀疏响应约束系统性地统一嵌入到生成模型的网络结构与训练目标之中。

\textbf{我们的方法。}针对上述问题，本研究提出PriorFlow——在条件流匹配框架的四个层面系统融入“基因身份感知”机制。（1）在条件编码层，以GO功能关联、STRING蛋白互作和TRRUST转录调控三层先验构建异质扰动基因图，通过图注意力网络（GATv2）替代静态嵌入，使扰动条件表示携带基因间的功能相似性、物理互作和调控方向信息。（2）在表达融合层，以扰动条件为query、表达基因为key/value进行交叉注意力计算，替代简单拼接，使模型根据具体扰动主动检索受影响的基因程序。（3）在速度场输出层，设计TRRUST偏置的gene mask，对每个基因施加可学习的稀疏响应调制，使已知调控靶点获得优先关注。（4）在训练目标层，引入基因级信噪比（SNR）加权终端损失，使变化幅度大且跨细胞稳定的高可信响应基因在训练中占据更大权重，从根本上缓解无响应基因梯度淹没的问题。

\subsection{基于因果解耦的扰动表示学习与条件生成}

本研究的第二个维度聚焦于因果解耦与条件流匹配的深度融合，旨在学习具有因果语义的扰动表示并实现高质量的条件生成，从而提升模型对扰动机制的解释能力和跨条件泛化性能。

\textbf{问题背景。}即便一个扰动预测模型能够准确地预测出哪些基因在扰动后发生了表达变化（即解决了研究点一的”预测精度”问题），一个更深层的挑战依然存在：预测到的基因表达变化中，哪些是扰动基因直接因果调控的结果，哪些仅仅是与扰动基因共表达的间接关联？这一区分的生物学意义重大——如果基因B的变化仅仅是因为它与真正响应基因A共表达，而非被扰动基因直接调控，那么针对基因B的后续干预将是无效的。然而，单细胞扰动数据本身仅提供扰动前后的表达快照，并不直接标注基因间的因果调控方向。更棘手的是，细胞类型、细胞周期、代谢状态等”背景”因素占据了表达变异的主导部分，扰动诱导的真实因果变化通常稀疏且微弱，极易被背景波动所掩盖——模型可能将两个细胞间的背景差异误判为扰动效应。

\textbf{现有方法的局限。}当前的主流方法（如CPA、CellFlow、GEARS等）本质上是学习从对照分布到扰动分布的统计映射，其输出是”扰动后表达谱的预测值”，而非”扰动因果效应的估计值”。PerturbedVAE\cite{jiang2025what}揭示了一个关键现象——扰动抑制假说（Perturbation Suppression Hypothesis）：模型容易被占主导的背景表达模式吸引而忽略稀疏的扰动特异性信号，通用预训练表示在扰动标签的线性可解码性上甚至弱于PCA基线。这一发现表明，若不显式地将背景信号与因果信号分离，模型学到的”扰动响应”中将不可避免地混杂大量背景噪声。LCD\cite{lorch2025latent}将基因表达建模为由随机微分方程控制的稳态扩散过程，从稳态快照中推断基因调控动态，但其因果效应提取方法（CLIPR）依赖于线性近似假设。总体而言，现有方法尚未在统一的生成框架中同时实现”背景与因果信号的显式解耦”与”基于解耦表示的高质量条件生成”——前者是确保扰动响应具有因果语义的前提，后者是使因果表示服务于实际预测任务的关键。

\textbf{我们的方法。}针对上述问题，本研究提出CausalFlow——将因果解耦表示学习与条件流匹配生成深度融合。首先，将基因表达的潜在表示显式分解为扰动不变变量$z_\iota$（捕捉细胞类型、细胞周期等背景转录程序）和扰动响应变量$z_\nu$（捕捉扰动诱导的特异性变化），其中$z_\nu$采用自回归结构以编码因果图中的拓扑排序关系。其次，引入对比对齐机制——强制同一细胞扰动前后的扰动不变表示保持一致，迫使所有扰动诱导的变化进入$z_\nu$，从而实现背景与因果信号的显式分离。随后，以解耦后的因果表示$[z_\iota; z_\nu]$作为条件输入流匹配生成模块，实现从基分布到扰动后转录组分布的条件传输。最后，借鉴CLIPR\cite{lorch2025latent}的思想并推广至非线性因果生成模型，利用多扰动条件下的响应矩阵提取基因间因果效应，并通过与TRRUST等已知调控关系数据库的比对进行系统验证。

\subsection{预期目标}

在方法层面，本课题拟构建一套完整的条件流匹配扰动预测方法体系，涵盖两个互补的维度：基于先验约束的稀疏响应扰动预测方法，将生物学先验知识与深度生成模型有机融合，重点提升关键响应基因、扰动效应方向和响应幅度的预测能力；因果解耦的条件流匹配框架，实现扰动因果表示的学习与高质量条件生成。目标是在预测精度、扰动响应识别能力和因果可解释性上全面超越现有方法。

在理论层面，拟从三个方向展开探索。其一，建立条件流匹配与生物学先验知识融合的理论框架，分析知识引导对稀疏扰动响应识别、差异表达基因恢复和生成质量的影响机理。其二，发展因果解耦与流匹配协同优化的理论分析，厘清因果表示可识别性与生成质量之间的内在关系。其三，探索从稳态单细胞快照数据中推断因果调控机制的理论条件与限制，为相关方法的适用性提供判据。

在工程层面，计划开发一套开源的单细胞扰动预测工具包，涵盖数据预处理、模型训练、预测生成和因果分析等完整流程，降低社区使用门槛。同时力争在国际顶级期刊或会议上发表1-2篇高水平学术论文。

\clearpage
\section{拟采用的研究方法、技术路线、实验方案及其可行性分析}

\subsection{条件流匹配算法基础}

条件流匹配（Conditional Flow Matching, CFM）是本课题的核心技术基础。流匹配的基本思想是通过学习从简单基分布$p_0$到复杂目标分布$p_1$的连续归一化流来生成数据。给定目标数据分布$p_1(x)$，条件流匹配通过定义条件概率路径$p_t(x|x_1)$来构建生成模型，其中$t \in [0,1]$，边界条件为$p_0(x|x_1) = p_0(x)$和$p_1(x|x_1) \approx \delta_{x_1}(x)$。令$u_t(x|x_1)$为生成条件概率路径的条件向量场，流匹配的训练目标为：
\begin{equation}
\mathcal{L}_{\text{CFM}}(\theta) = \mathbb{E}_{t \sim \mathcal{U}(0,1), x_1 \sim p_1, x \sim p_t(\cdot|x_1)}\|v_t(x; \theta) - u_t(x|x_1)\|^2
\end{equation}
其中$v_t(x; \theta)$为由神经网络参数化的学习向量场。在推理阶段，给定初始样本$x_0 \sim p_0$，通过求解常微分方程（ODE）$\frac{d}{dt}x(t) = v_t(x(t); \theta)$生成目标样本。

条件流匹配相比扩散模型具有以下优势：（1）训练目标直接回归向量场，避免了得分匹配中的复杂优化；（2）采样路径是确定性的，采样效率和稳定性更好；（3）支持灵活的基分布和条件路径设计，易于整合各类先验约束。这些特性使得条件流匹配特别适合作为单细胞扰动预测的生成骨架。

\subsection{基于先验增强的稀疏响应扰动预测}

\subsubsection{总体算法框架}

本方法以最优传输条件流匹配（OT-CFM）为生成骨架，将单细胞扰动预测表述为从对照分布到扰动分布的条件连续传输问题。给定对照细胞$x_0\in\mathbb{R}^d$、真实扰动细胞$x_1\in\mathbb{R}^d$和扰动条件$c$（$d$为表达基因数），模型学习条件向量场$v_\theta(x_t,t|c)$，并通过求解ODE从对照细胞出发积分得到目标扰动状态的预测。

然而，高维转录组空间中存在一个核心挑战：在数千个表达基因中，真正对特定扰动产生显著响应的基因通常仅占少数（通常是数十到数百个差异表达基因），大量无响应基因的梯度信号会淹没响应基因的监督信息。现有方法（如CPA、CellFlow、scDFM）普遍使用one-hot编码或静态可学习嵌入来表示扰动条件，忽视了基因之间客观存在的功能关联、蛋白互作和转录调控关系——而这些先验信息恰恰是理解扰动信号如何沿生物网络传播的关键线索。此外，标准流匹配损失对所有基因维度等权求和，无法区分高信噪比响应基因与随机噪声基因。

针对上述问题，本方法将”扰动响应通常集中于少数关键基因或通路”的生物规律系统性地融入条件流匹配的四个层面：（1）条件编码层：构建融合GO功能关联、STRING蛋白互作和TRRUST转录调控的三层扰动基因先验图，通过图注意力网络获得结构增强的扰动条件表示；（2）表达融合层：以扰动条件表示为query、表达基因为key/value进行交叉注意力计算，显式建模扰动基因与表达基因之间的调控交互；（3）速度场输出层：从条件嵌入预测基因级稀疏掩码，并以TRRUST调控关系为偏置，在基因粒度上约束速度场的学习方向；（4）训练目标层：设计SNR（信噪比）加权的终端MSE损失，按基因级效应大小与方差之比分配权重，强化高响应基因的监督信号。以下各小节逐一详述这四个层面的设计与实现。

\subsubsection{先验扰动图构建与图增强条件编码}

\textbf{扰动标识统一与条件表示。} 首先将不同数据集中的扰动标识（药物名、sgRNA靶基因等）统一映射到基因符号空间，保证扰动标签、表达基因列表和先验图节点三者之间的一致对应。对于组合扰动，一个条件$c$被表示为扰动基因集合$\{g_1,\ldots,g_k\}$；对于单基因扰动，$k=1$。模型为每个实际出现的扰动基因维护一个可学习的初始嵌入向量$e_g^{\mathrm{pert}}\in\mathbb{R}^{d_{\mathrm{gnn}}}$（默认$d_{\mathrm{gnn}}=128$），以正态分布$\mathcal{N}(0,0.02)$随机初始化。这些嵌入随后在扰动基因先验图上通过图神经网络进行结构增强。

\textbf{三层先验图的构建。} 先验图$\mathcal{G}=(\mathcal{V},\mathcal{E})$仅在训练集中实际出现的扰动基因集合上构建，节点数通常为数百到数千量级（取决于数据集），远小于全基因组规模，既保留了关键的生物学先验，又避免了在全基因组的无关节点上过度平滑。图中融合了三种互补的生物学关系类型：

\textit{GO功能关联。} 基于基因本体（Gene Ontology）注释计算基因间的功能语义相似性。对于每一对基因，计算其共享GO term的信息含量（Information Content, IC）之和与全部GO term的IC之和的比值$s_{ij}=IC_{\mathrm{shared}}/IC_{\mathrm{union}}$，其中每个GO term的IC定义为$\mathrm{IC}(t)=\log\frac{N_{\mathrm{total}}+1}{N_t+1}$，$N_t$为注释到该term的基因数，$N_{\mathrm{total}}$为总基因数。为控制图的稀疏性，对每个目标基因仅保留相似度最高的top-k个源基因（默认$k=20$），并可通过幂次缩放$w_{ij}=s_{ij}^{p}$（默认$p=1.0$）调节边权分布。

\textit{STRING蛋白互作。} 从STRING数据库获取实验验证或数据库证据支持的蛋白-蛋白物理相互作用。STRING数据以三元组（regulator, target, weight）形式存储于parquet文件中，其中weight为综合得分。仅保留两端均为扰动基因的边，并通过ppi\_weight\_scale系数（默认0.5）调节PPI边相对GO边的权重。当数据集中扰动基因数量较少导致GO边不足时（代码中阈值为100条），自动引入STRING边作为补充。

\textit{TRRUST转录调控。} 从TRRUST v2数据库\cite{han2018trrust}导入经文献人工整理的转录因子（TF）到靶基因的有向调控关系。与前两种边类型不同，TRRUST边连接的是扰动基因（作为TF）与表达基因（作为靶基因），因此不直接参与扰动基因图的消息传递，而是作为后续交叉注意力偏置和gene mask偏置的先验约束。在数据结构上，预处理阶段遍历所有扰动基因，对每个扰动基因查找其在TRRUST中的靶基因集合，构建一个二值矩阵$M_{\mathrm{trrust}}\in\{0,1\}^{N_{\mathrm{pert}}\times d}$，其中$M_{\mathrm{trrust}}[i,j]=1$当且仅当扰动基因$i$是表达基因$j$的已知转录调控因子。运行时根据当前扰动条件涉及的扰动基因索引，从$M_{\mathrm{trrust}}$中取出对应行并做max-pooling聚合，得到该条件下所有表达基因的二值靶标掩码$m_{\mathrm{trrust}}\in\{0,1\}^{d}$。

\textbf{边权归一化。} 为缓解高度数节点的消息主导效应，对GO和STRING合并后的所有边进行目标节点度归一化：
\begin{equation}
\tilde{w}_{ij}=\frac{w_{ij}}{\sum_{\ell\in\mathcal{N}_{\mathrm{in}}(j)}w_{\ell j}+\epsilon},
\label{eq:degree_norm}
\end{equation}
其中$\mathcal{N}_{\mathrm{in}}(j)$为指向节点$j$的所有源节点集合，$\epsilon=10^{-8}$防止除零。

\textbf{图消息传递网络。} 模型提供了两种图编码器变体，通过配置切换。

\textit{基础版本PerturbationGNN}采用固定边权的消息传递。记第$l$层节点特征为$H^{(l)}\in\mathbb{R}^{N_{\mathrm{pert}}\times d_{\mathrm{gnn}}}$，消息传递过程为：
\begin{equation}
M_j^{(l)}=\sum_{i\in\mathcal{N}_{\mathrm{in}}(j)}\tilde{w}_{ij}\cdot H_i^{(l)},\qquad
H^{(l+1)}=\mathrm{LayerNorm}\left(H^{(l)}+\mathrm{MLP}\left(M^{(l)}\right)\right),
\label{eq:gnn_basic}
\end{equation}
其中MLP为两层线性变换（$d_{\mathrm{gnn}}\to d_{\mathrm{gnn}}\to d_{\mathrm{gnn}}$），中间使用silu激活，最后一层无激活。默认堆叠$L=2$层，每层后应用LayerNorm残差连接。

\textit{增强版本EnhancedPerturbationGNN}在基础版本之上引入了三个关键改进：（i）用多头GATv2注意力替代固定权重的消息传递，注意力分数由源节点query与目标节点key相加后经LeakyReLU得到，边权以对数空间偏置$\log(\tilde{w}_{ij})$的形式注入注意力计算；(ii)引入一个可学习的虚拟节点$\mathbf{vn}\in\mathbb{R}^{1\times d_{\mathrm{gnn}}}$作为全局通信枢纽——虚拟节点通过注意力机制聚合所有真实节点的信息，随后真实节点通过门控机制$g=\sigma(\mathrm{score}(h_i,\mathbf{vn}))$有选择地融合全局信息；(iii)将层数增加到$L=4$，每层采用Pre-norm残差结构，包含GATv2消息传递、门控全局信息融合和两层FFN（$d_{\mathrm{gnn}}\to 2d_{\mathrm{gnn}}\to d_{\mathrm{gnn}}$，gelu激活），默认使用4个注意力头。虚拟节点的设计灵感受Exphormer启发，但以可学习的虚拟节点替代了随机扩展图，提供了更具可解释性的全局通信路径。训练时对GNN层施加独立的graph\_dropout正则化。

\textbf{集合注意力池化。} 经过图编码后，样本$b$的$k$个扰动基因对应的增强嵌入$\{h_{b,1}^{\mathrm{gnn}},\ldots,h_{b,k}^{\mathrm{gnn}}\}$通过集合注意力池化（TokenAttentionPooling）聚合为固定维度的条件嵌入$z_c^{(b)}\in\mathbb{R}^{d_c}$（Norman additive实验$d_c=256$，Replogle LOCO实验$d_c=512$）：
\begin{equation}
z_c^{(b)}=\mathrm{AttentionPool}\left(\{h_{b,1}^{\mathrm{gnn}},\ldots,h_{b,k}^{\mathrm{gnn}}\}\right),
\label{eq:set_pool}
\end{equation}
该池化机制通过可学习的注意力token自适应地聚合变长扰动基因集合，对padding位置施加掩码以排除空位干扰。池化后的嵌入可选通过额外的MLP层（layers\_after\_pool）进一步变换，并支持确定性或随机（重参数化高斯采样）两种编码模式。

\subsubsection{扰动条件到表达基因的交叉注意力融合}

得到扰动条件嵌入$z_c$后，模型需要将其与当前细胞状态$x_t$融合，以生成条件感知的表示。与常见的将条件嵌入简单拼接到细胞状态或时间编码上的做法不同，本方法在基因粒度上建模”扰动条件如何影响每个表达基因”的交互关系。

\textbf{表达基因的身份-数值双通道表示。} 对于路径时刻$t$的细胞状态$x_t\in\mathbb{R}^{d}$，模型首先将每个表达基因$i$的表达值通过线性投影映射到注意力维度空间，并叠加一个可学习的基因身份嵌入，形成”表达量（what）+ 基因身份（who）”的双通道表示：
\begin{equation}
h_i^{\mathrm{val}}=W_vx_{t,i}\in\mathbb{R}^{d_{\mathrm{attn}}},\qquad
h_i^{\mathrm{id}}=e_i^{\mathrm{gene}}\in\mathbb{R}^{d_{\mathrm{attn}}},\qquad
H_x=[h_1^{\mathrm{val}}+h_1^{\mathrm{id}},\;\ldots,\;h_d^{\mathrm{val}}+h_d^{\mathrm{id}}]\in\mathbb{R}^{d\times d_{\mathrm{attn}}},
\label{eq:gene_repr}
\end{equation}
其中$W_v\in\mathbb{R}^{d_{\mathrm{attn}}\times 1}$将标量表达值投影到$d_{\mathrm{attn}}$维（Norman additive实验16维，Replogle LOCO实验64维），基因身份嵌入$E^{\mathrm{gene}}\in\mathbb{R}^{d\times d_{\mathrm{attn}}}$以$\mathcal{N}(0,0.02)$初始化并在训练中学习。这种设计使模型在注意力计算中既能感知表达量的高低，又能通过身份嵌入区分不同基因——即使两个基因在当前细胞中表达量相同，模型也能通过身份嵌入将它们区别对待。

\textbf{条件到基因的交叉注意力。} 模型以扰动条件嵌入$z_c$为query、以表达基因表示$H_x$为key和value，执行交叉注意力计算：
\begin{equation}
Q=W_qz_c\in\mathbb{R}^{1\times d_{\mathrm{attn}}},\quad
K=W_kH_x,\;V=W_vH_x\in\mathbb{R}^{d\times d_{\mathrm{attn}}},
\label{eq:cross_attn_qkv}
\end{equation}
\begin{equation}
h_{\mathrm{attn}}=\mathrm{MultiHeadAttn}(Q,\;K,\;V)\in\mathbb{R}^{1\times d_{\mathrm{attn}}},
\label{eq:cross_attn}
\end{equation}
其中多头注意力使用4个头，实验中均为1层。当配置多层交叉注意力时（cross\_attn\_layers > 1），每层的注意力输出通过广播残差连接回基因表示$H_x\leftarrow H_x+h_{\mathrm{attn}}$，使基因表示随注意力层逐步精化——底层关注广泛的响应基因模式，高层聚焦于条件特异的关键调控靶点。交叉注意力输出经LayerNorm归一化和线性投影（fusion\_proj, $d_{\mathrm{attn}}\to d_{\mathrm{hidden}}$，其中$d_{\mathrm{hidden}}=512$）后，与时间编码拼接。时间编码由正弦位置编码（频率数1024，最大周期10000）经三层MLP（512→512→512，silu激活）得到。拼接后的向量经三层解码器MLP（1024→1024→1024，silu激活，默认dropout rate 0.1）和最终线性层映射为基因级速度场。可选地，在交叉注意力之前对表达基因特征施加基因自注意力（gene\_self\_attn\_layers），但实验中均设为0（$O(d^2)$计算量过大）。

\textbf{TRRUST目标基因的注意力特征标记（方案B）。} 除上述交叉注意力机制外，模型可选启用TRRUST注意力偏置：为TRRUST已知的靶基因在KV空间中附加可学习的特征标记$e_{\mathrm{trrust}}\in\mathbb{R}^{1\times d_{\mathrm{attn}}}$：
\begin{equation}
H_x\leftarrow H_x + m_{\mathrm{trrust}}\otimes e_{\mathrm{trrust}},
\label{eq:trrust_attn_bias}
\end{equation}
其中$m_{\mathrm{trrust}}\in\{0,1\}^{d}$为当前扰动条件下TRRUST靶基因的二值掩码（由$M_{\mathrm{trrust}}$矩阵按扰动基因索引查找并取max-pooling聚合得到）。这一设计使已知的转录调控靶点在交叉注意力的KV空间中具备可区分的调控信号，帮助注意力机制更精确地检索受当前扰动调控的关键基因。

\subsubsection{TRRUST偏置的基因级稀疏速度场调制}

扰动响应具有天然的稀疏性——一个扰动通常只显著改变少数基因的表达。为将这一归纳偏置注入速度场学习，模型从扰动条件嵌入出发，预测一个基因级的重要性掩码，直接调制速度场的输出。

\textbf{条件感知的基因掩码预测。} 在速度场的输出端，模型通过一个线性头从条件嵌入预测每个表达基因的重要性分数：
\begin{equation}
m_c^{\mathrm{raw}}=W_mz_c+b_m\in\mathbb{R}^{d},
\label{eq:mask_raw}
\end{equation}
其中$W_m\in\mathbb{R}^{d\times d_c}$。当启用TRRUST掩码偏置（方案A）时，将TRRUST靶基因的二值掩码以可学习强度注入：
\begin{equation}
m_c^{\mathrm{raw}}\leftarrow m_c^{\mathrm{raw}} + \alpha\cdot m_{\mathrm{trrust}},
\label{eq:mask_trrust}
\end{equation}
其中$\alpha\in\mathbb{R}$为可学习标量参数，初始化为1.0，使模型在训练初期即倾向于关注已知调控靶点，随后可自适应调节偏置强度。

\textbf{均值归一化与速度场调制。} 原始掩码经sigmoid激活后，进行均值归一化以确保掩码仅重新分配基因维度间的相对重要性，而非全局缩放速度场幅值（后者会导致模型通过缩小输出幅度来”欺骗”损失函数）：
\begin{equation}
m_c=\frac{\sigma(m_c^{\mathrm{raw}})}{\frac{1}{d}\sum_{i=1}^{d}\sigma(m_c^{\mathrm{raw}})_i+\epsilon},
\label{eq:mask_norm}
\end{equation}
最终的稀疏调制速度场为原始速度场与归一化掩码的逐元素乘积：
\begin{equation}
v_\theta^{\mathrm{mask}}(x_t,t|c)=v_\theta(x_t,t|c)\odot m_c.
\label{eq:masked_velocity}
\end{equation}

该设计的核心思想在于：gene mask从条件嵌入中学习”该扰动倾向于影响哪些基因”，本质上是将全局条件信息解码为基因粒度的调控重要性分布。与在损失函数中间接施加稀疏约束（如L1正则）不同，在速度场输出端直接masking的方式将稀疏先验结构化地嵌入模型推理过程——模型不需要额外学习”哪些基因不应响应”，因为无响应基因的输出天然被mask抑制。此外，TRRUST偏置提供了一条从文献验证的调控关系到模型参数的捷径，降低了模型从零开始学习调控靶点的难度。

\subsubsection{SNR加权终端损失与多目标训练}

\textbf{流匹配基础损失。} 训练时从对照细胞$x_0$、扰动细胞$x_1$和条件$c$构造条件概率路径$x_t=(1-t)x_0+tx_1$（线性插值路径），并通过Sinkhorn最优传输匹配形成训练对，以减小条件分布之间的传输代价。基础流匹配损失为标准MSE：
\begin{equation}
\mathcal{L}_{\mathrm{FM}}=\mathbb{E}_{t,x_0,x_1,c}\left\|v_\theta^{\mathrm{mask}}(x_t,t|c)-u_t(x_t|x_0,x_1)\right\|^2,
\label{eq:loss_fm}
\end{equation}
其中$u_t=x_1-x_0$为线性路径的真实向量场。

\textbf{SNR加权终端损失。} 基础流匹配损失对全部$d$个基因等权平均。在单细胞扰动场景中，大量基因在扰动前后表达水平变化很小或纯属随机波动，这部分基因的梯度信号会淹没少数真正响应基因的监督信息。为缓解此问题，本文提出基因级信噪比（Signal-to-Noise Ratio）加权方案。

对于给定条件$c$，定义基因$g$的扰动效应$\delta_g=x_{1,g}-x_{0,g}$，在同条件的所有训练样本上计算其条件均值（效应大小）和条件标准差（效应稳定性）：
\begin{equation}
\mu_{c,g}=\mathbb{E}_{(x_0,x_1)\sim p(\cdot|c)}[\delta_g],\qquad
\sigma_{c,g}=\sqrt{\mathbb{E}_{(x_0,x_1)\sim p(\cdot|c)}[(\delta_g-\mu_{c,g})^2]+\epsilon},
\label{eq:snr_stats}
\end{equation}
其中$\epsilon=10^{-8}$。$\mu_{c,g}$反映扰动在基因$g$上引起的平均效应大小，$\sigma_{c,g}$反映该效应在同条件不同细胞间的一致性。若一个基因效应大且稳定（$|\mu|$大、$\sigma$小），则其为高置信响应基因；若效应小或波动大，则为低信噪比基因。

基因$g$的原始SNR权重定义为类Cohen's d效应量：
\begin{equation}
s_{c,g}=\frac{|\mu_{c,g}|}{\sigma_{c,g}+\epsilon}.
\label{eq:snr_raw}
\end{equation}
为保持损失量级稳定，对SNR权重进行均值归一化：
\begin{equation}
\omega_{c,g}=\frac{s_{c,g}}{\frac{1}{d}\sum_{g'=1}^{d}s_{c,g'}+\epsilon},
\label{eq:snr_norm}
\end{equation}
使得$\mathbb{E}_g[\omega_{c,g}]\approx 1$，整体损失量级不会因子基因权重重新分配而改变。此外，对权重施加裁剪$\omega_{c,g}\leftarrow\min(\omega_{c,g},w_{\max})$（默认$w_{\max}=4.0$），防止极少数超高分基因主导梯度。可选地对top-k（默认前5\%或至少20个）最大$|\mu_{c,g}|$基因施加额外加权系数。

SNR加权的终端MSE损失定义如下：令端点预测为$\hat{x}_1=x_t+(1-t)v_\theta^{\mathrm{mask}}(x_t,t|c)$，则
\begin{equation}
\mathcal{L}_{\mathrm{SNR}}=\mathbb{E}_{t,x_0,x_1,c}\left[t^{p}\sum_{g=1}^{d}\omega_{c,g}\left(\hat{x}_{1,g}-x_{1,g}\right)^2\right],
\label{eq:loss_snr}
\end{equation}
其中$t^{p}$为时间门控因子（默认$p=2$），使终端损失在$t\to 1$（接近目标分布）时权重更大，避免在流匹配早期阶段噪声较大的预测上施加过强的端点约束。SNR权重$\omega_{c,g}$在训练过程中不参与梯度反向传播（视为常数），仅作为损失加权系数使用。

\textbf{综合训练目标。} 完整训练损失为流匹配基础损失与SNR加权终端损失的组合：
\begin{equation}
\mathcal{L}=
\mathcal{L}_{\mathrm{FM}}
+\lambda_{\mathrm{SNR}}\mathcal{L}_{\mathrm{SNR}},
\label{eq:loss_total}
\end{equation}
其中$\lambda_{\mathrm{SNR}}$为SNR损失权重（Norman additive实验设为1.0，Replogle LOCO实验设为2.0）。

\textbf{推理过程。} 推理时取对照细胞$x_0$，将扰动条件$c$输入模型获得条件嵌入（使用确定性模式，即直接输出均值），在时间区间$[0,1]$上以Tsit5自适应步长求解器求解ODE初值问题$dx/dt=v_\theta^{\mathrm{mask}}(x,t|c),\;x(0)=x_0$，得到$x(1)$即为预测的扰动表达谱。求解器使用PID步长控制器（rtol=$10^{-5}$, atol=$10^{-5}$）以保证积分精度。

\subsubsection{实验设置与消融验证}

实验上设置两类具有代表性的评估场景。第一类为Norman additive split，测试部分双基因组合扰动，训练集中保留相关单基因扰动，用于评估模型对组合扰动响应中关键差异表达基因和效应方向的预测能力。训练配置上使用EnhancedPerturbationGNN（GATv2, $d_{\mathrm{gnn}}=128$, $L=4$, 4头），条件嵌入维度256，基因注意力维度16，交叉注意力4头1层，无基因自注意力，流噪声0.0，SNR权重1.0，训练10000轮，学习率$5\times10^{-4}$，批量大小256。第二类为Replogle LOCO split，留出hepg2细胞系中的部分扰动响应，用于评估模型在不同细胞背景下预测扰动效应方向、幅度和差异表达基因的能力。训练配置上同样使用EnhancedPerturbationGNN（$d_{\mathrm{gnn}}=128$, $L=4$, 4头），条件嵌入维度512，基因注意力维度64，交叉注意力4头1层，无基因自注意力，流噪声0.1，SNR权重2.0，最优传输匹配间隔20轮，训练20000轮，学习率$5\times10^{-4}$，批量大小256。两个实验均启用TRRUST的gene mask偏置（方案A）和交叉注意力偏置（方案B）。对比方法包括GEARS、CellFlow、scDFM、TxPert和CPA，覆盖图神经网络、流匹配、扩散模型、Transformer/GNN和自编码器等不同建模范式。

评价指标包括MSE、Pearson delta、top-k delta相关、DEG集合上的PCC/R2/EV、DEG overlap F1以及方向一致性等。所有扰动效应相关指标均按扰动条件分别计算后再取平均，避免不同扰动方向在全局平均中相互抵消。消融实验将重点评估GO+STRING扰动图、交叉注意力融合、gene mask和SNR加权损失四个组件的贡献。

\begin{table}[ht]
\centering
\small
\caption{基于先验约束的扰动预测实验数据划分设置}
\begin{tabular}{|p{3.2cm}|p{3.0cm}|p{4.0cm}|p{4.0cm}|}
\hline
\textbf{数据集与任务} & \textbf{评估场景} & \textbf{训练集设置} & \textbf{测试集设置} \\
\hline
Norman additive & 组合扰动响应评估 & 保留全部单基因扰动，并使用部分双基因组合扰动 & 留出部分双基因组合扰动，用于评估差异表达基因和扰动方向预测能力 \\
\hline
Replogle LOCO & 跨细胞背景响应评估 & 使用非留出细胞系扰动响应，并保留留出细胞系对照细胞 & 留出hepg2细胞系中的部分扰动响应，用于评估效应方向、幅度和差异表达基因预测能力 \\
\hline
\end{tabular}
\end{table}

\begin{table}[ht]
\centering
\small
\caption{Norman additive组合扰动预测主结果表}
\begin{tabular}{|p{2.7cm}|c|c|c|c|c|}
\hline
\textbf{方法} & \textbf{MSE$\downarrow$} & \textbf{Pearson $\Delta\uparrow$} & \textbf{Top-20 $\Delta\uparrow$} & \textbf{DEG PCC$\uparrow$} & \textbf{DEG F1$\uparrow$} \\
\hline
CPA & -- & -- & -- & -- & -- \\
\hline
GEARS & -- & -- & -- & -- & -- \\
\hline
CellFlow & -- & -- & -- & -- & -- \\
\hline
scDFM & -- & -- & -- & -- & -- \\
\hline
TxPert & -- & -- & -- & -- & -- \\
\hline
PriorFlow（本文） & -- & -- & -- & -- & -- \\
\hline
\end{tabular}
\end{table}

\begin{table}[ht]
\centering
\small
\caption{Replogle LOCO跨细胞系扰动预测主结果表}
\begin{tabular}{|p{2.7cm}|c|c|c|c|c|}
\hline
\textbf{方法} & \textbf{MSE$\downarrow$} & \textbf{Pearson $\Delta\uparrow$} & \textbf{Top-20 $\Delta\uparrow$} & \textbf{DEG PCC$\uparrow$} & \textbf{DEG F1$\uparrow$} \\
\hline
CPA & -- & -- & -- & -- & -- \\
\hline
GEARS & -- & -- & -- & -- & -- \\
\hline
CellFlow & -- & -- & -- & -- & -- \\
\hline
scDFM & -- & -- & -- & -- & -- \\
\hline
TxPert & -- & -- & -- & -- & -- \\
\hline
PriorFlow（本文） & -- & -- & -- & -- & -- \\
\hline
\end{tabular}
\end{table}

\begin{table}[ht]
\centering
\small
\caption{PriorFlow关键模块消融实验表}
\begin{tabular}{|p{4.2cm}|c|c|c|c|}
\hline
\textbf{模型变体} & \textbf{Pearson $\Delta\uparrow$} & \textbf{Top-20 $\Delta\uparrow$} & \textbf{DEG PCC$\uparrow$} & \textbf{DEG F1$\uparrow$} \\
\hline
完整PriorFlow & -- & -- & -- & -- \\
\hline
w/o GO+STRING扰动图 & -- & -- & -- & -- \\
\hline
w/o 条件到基因交叉注意力 & -- & -- & -- & -- \\
\hline
w/o gene mask & -- & -- & -- & -- \\
\hline
w/o SNR加权损失 & -- & -- & -- & -- \\
\hline
\end{tabular}
\end{table}

\subsection{基于因果解耦的扰动表示学习与条件生成}

\subsubsection{潜在因果表示建模}

基于第3.2节描述的理论框架，构建包含扰动不变变量$z_\iota$和扰动响应变量$z_\nu$的潜在因果生成模型。采用变分推断方法进行模型训练，使用编码器网络$q_\phi(z_\iota, z_\nu | x, u)$近似后验分布，解码器网络$p_\psi(x | z_\iota, z_\nu)$重构基因表达。变分后验采用结构化分解：
\begin{align}
q(z_\iota | x) &= \mathcal{N}(\mu'_\iota(x), \text{diag}(\sigma'_\iota(x))) \\
q(z_\nu | x, u) &= \prod_{i=1}^{d_\nu} q(z_{\nu,i} | z_{\nu,<i}, x, u) = \prod_{i=1}^{d_\nu} \mathcal{N}(\mu'_{\nu,i}, \sigma'_{\nu,i})
\end{align}
其中$q(z_\nu | x, u)$采用自回归结构以捕捉DAG编码的因果依赖关系。

\subsubsection{对比对齐训练}

实施对比对齐损失，通过强制扰动样本和对照样本的$z_\iota$一致性实现因果解耦。训练采用联合优化策略：
\begin{equation}
\mathcal{L}_{\text{total}} = \mathcal{L}_{\text{rec}} + \beta_\nu \mathcal{L}_{\text{KL-}\nu} + \beta_\iota \mathcal{L}_{\text{KL-}\iota} + \alpha \mathcal{L}_{\text{contrast}}
\end{equation}
其中$\mathcal{L}_{\text{rec}} = \|x - \hat{x}\|^2$为重构损失，$\mathcal{L}_{\text{KL-}\nu}$和$\mathcal{L}_{\text{KL-}\iota}$分别为扰动响应变量和扰动不变变量的KL正则化项，$\beta_\nu, \beta_\iota, \alpha$为平衡系数。训练过程中采用退火策略，逐步增加KL项的权重。

\subsubsection{因果条件流匹配生成}

在获得解耦的因果表示后，构建以因果表示为条件的流匹配生成模型。向量场网络接受拼接的因果表示$h = [z_\iota; z_\nu]$作为条件，学习从噪声到基因表达分布的传输映射。为实现端到端的训练，探索两种方案：（a）两阶段训练——先训练因果解耦模块，再基于固定的因果表示训练流匹配模块；（b）端到端联合训练——同时优化因果解耦和流匹配目标，使因果表示学习与条件生成相互促进。

\subsubsection{因果效应提取与验证}

将CLIPR方法推广至非线性因果生成模型。对于每个扰动条件$q$，计算初始扰动响应和极限扰动响应，并利用多个扰动条件下的响应矩阵求解基因间因果效应。对提取的因果结构进行系统的生物学验证：将预测的因果边与已知调控关系数据库（如TRRUST、RegNetwork）进行比对；对因果效应最强的上游调控基因进行通路富集分析；验证预测因果边的方向性是否与已知调控方向一致。

\subsection{实验方案}

（1）数据集：
\begin{itemize}
    \item \textbf{Norman et al. (2019)}：K562细胞系的CRISPRa Perturb-seq数据，包含约96,804个细胞，覆盖105个单基因扰动和131个双基因扰动组合，是评估组合扰动响应预测和差异表达基因恢复能力的标准基准。
    \item \textbf{Replogle et al. (2022)}：K562细胞系的CRISPRi Perturb-seq数据，包含约209,462个细胞，覆盖约1000个基因的全基因组扰动筛选，是评估大规模扰动预测能力的重要数据集。
    \item \textbf{Wessels et al. (2023)}：THP-1细胞系的Cas13 RNA Perturb-seq数据，包含约23,722个细胞，覆盖25个基因的敲低扰动。
\end{itemize}

（2）评估指标：
\begin{itemize}
    \item \textbf{最大均值差异（MMD）}：衡量生成分布与真实分布之间的差异，使用平方指数核$k(y, y') = \exp(-\|y-y'\|^2 / 2d\gamma^2)$，在10个对数间隔的长度尺度上取平均。
    \item \textbf{均方根误差（RMSE）}：衡量预测均值与观测均值之间的偏差。
    \item \textbf{皮尔逊相关系数（PCC）与$R^2$}：衡量预测表达与真实表达的相关性和解释方差比例。
    \item \textbf{差异表达基因恢复率}：评估模型对top-20差异表达基因的恢复能力。
    \item \textbf{因果结构恢复精度}：对于合成数据，使用AUROC评估预测因果边与真实因果网络的匹配度；使用F1评估top-$k$因果效应的分类准确性。
\end{itemize}

（3）基线方法：
\begin{itemize}
    \item \textbf{CPA}\cite{lotfollahi2023contextualization}：基于组合式VAE的扰动预测方法。
    \item \textbf{GEARS}\cite{roohani2024gears}：基于图神经网络的扰动预测方法。
    \item \textbf{CellFlow}\cite{huguet2024cellflow}：基于流匹配的扰动预测方法。
    \item \textbf{PerturbDiff}\cite{roman2024perturbdiff}：基于扩散模型的扰动预测方法。
    \item \textbf{SAMS-VAE / sVAE+}\cite{bereket2023sams, lopez2022svae}：基于稀疏结构/机制偏移的潜在变量模型。
    \item \textbf{Discrepancy-VAE}\cite{zhang2023identifiability}：基于分布差异的因果表示学习方法。
    \item \textbf{LCD}\cite{lorch2025latent}：基于潜在因果扩散的扰动预测方法。
\end{itemize}

（4）实验设计：
\begin{itemize}
    \item \textbf{单基因扰动预测}（分布内）：训练集和测试集的扰动类型相同，评估模型的拟合能力和生成质量。
    \item \textbf{双基因组合扰动预测}（分布外）：训练集仅包含对照和单基因扰动数据，测试集包含双基因扰动组合，评估模型对组合扰动效应、关键响应基因和方向一致性的刻画能力。
    \item \textbf{消融实验}：逐步移除GO+STRING扰动图、gene mask、SNR加权损失、对比对齐损失和因果结构先验等组件，量化各组件的贡献。
    \item \textbf{因果结构验证}：在合成数据（已知真实因果图）上评估因果效应提取的准确性；在真实数据上与已知调控关系数据库进行系统性比对。
\end{itemize}

\subsection{可行性分析}

条件流匹配框架在图像生成、音频合成等领域的成功应用已充分验证了其理论有效性，而因果表示学习的可识别性理论（如Zhang等人\cite{zhang2023identifiability}和Jiang等人\cite{jiang2025what}的近期工作）为将因果结构引入生成模型提供了严格的理论保障。LCD\cite{lorch2025latent}和PerturbedVAE\cite{jiang2025what}在单细胞扰动预测上的最新成果进一步表明，因果解耦与深度生成模型的融合具有切实的潜力。

数据方面，Norman等人（2019）、Replogle等人（2022）和Wessels等人（2023）发布的公开数据集覆盖了不同的扰动类型（CRISPRa、CRISPRi、Cas13）、细胞系和实验条件，为模型的训练和全面评估提供了充足的数据基础。技术实现上，PyTorch等深度学习框架和Scanpy、AnnData等单细胞分析工具均已高度成熟，条件流匹配和变分推断也有torchcfm、scvi-tools等高质量开源参考实现可供定制化开发。

计算资源方面，本研究所需的模型训练可在配备NVIDIA A100（40GB及以上显存）的GPU集群上完成。LCD的两阶段训练中，先验密度推断阶段虽需较大蒙特卡洛样本量（$M \approx 20,000$），但大温度$\tau$近似等加速策略已能有效控制计算开销，整体训练和推理成本在可控范围内。

从创新可行性来看，本课题的两个研究点均有清晰的技术路径和分阶段验证方案。基于先验约束的稀疏响应扰动预测可在现有流匹配代码基础上扩展实现，并通过GO+STRING扰动图、gene mask和SNR加权损失的消融实验逐步验证；因果解耦与流匹配的融合则可以先独立验证因果解耦模块的可识别性和流匹配模块的生成质量，再推进端到端联合优化，风险可控。

\clearpage
\section{已有科研基础与所需的科研条件}

\subsection{已有科研基础}

在文献调研方面，已完成了对单细胞扰动预测领域的系统梳理，深入阅读了流匹配、扩散模型、因果表示学习、基因调控网络建模等方向的核心文献80余篇，对该领域的发展脉络和前沿动态有了较为全面的把握。

方法储备上，已系统学习了条件流匹配与连续归一化流的理论及实现、变分自编码器与变分推断方法、因果表示学习与可识别性理论、图神经网络及其生物信息学应用，以及随机微分方程的数值求解方法。数据方面，已收集并完成了Norman等人（2019）、Replogle等人（2022）、Wessels等人（2023）等公开数据集的预处理和探索性分析。

在先期实践方面，已复现了CellFlow和PerturbDiff的核心算法，对条件流匹配在单细胞数据上的表现取得了初步实验认知；同时阅读并理解了LCD和PerturbedVAE的源代码框架，为后续开发工作奠定了基础。编程方面熟练掌握Python、PyTorch、Scanpy、AnnData等工具链，具备独立开展深度学习研究的能力。

\subsection{所需科研条件}

硬件方面，需要配备NVIDIA A100（40GB显存及以上）或同等级别GPU的计算服务器，用于大规模模型训练和超参数搜索；约500GB存储空间用于存放原始数据集、预处理数据和模型检查点。软件方面，需要PyTorch 2.0及以上版本、Scanpy、AnnData、torchcfm、scvi-tools等核心软件包，以及GSEApy、NetworkX等生物信息学分析工具。

学术资源方面，需要持续访问PubMed、bioRxiv、Google Scholar等学术数据库，以及ICLR、NeurIPS、ICML、ISMB、RECOMB等国际会议的最新论文集。同时，参加相关学术会议和研讨会、与国内外同行保持交流合作，对于及时了解领域前沿进展也是必要的。

\clearpage
\section{研究工作计划与进度安排}

\begin{table}[ht]
\centering
\small
\caption{研究工作计划与进度安排}
\begin{tabular}{|c|p{5.5cm}|p{4.5cm}|}
\hline
\textbf{时间段} & \textbf{工作内容} & \textbf{预期成果} \\
\hline
2026年6月--2026年8月 &
文献深入调研与开题报告完善；
公开数据集的收集、预处理和探索性分析；
条件流匹配和因果表示学习基础框架的搭建与初步实验 &
完成开题报告；
构建标准化基准数据集；
完成基线方法的复现与评估 \\
\hline
2026年9月--2026年12月 &
基于先验约束的稀疏响应扰动预测方法设计与实现；
知识图谱构建与知识感知编码器的开发；
gene mask与SNR加权损失等稀疏响应约束的设计与验证；
撰写研究论文初稿 &
完成第一个研究点的算法设计与实验验证；
形成一篇可投稿的研究论文 \\
\hline
2027年1月--2027年3月 &
因果解耦生成模型的构建与训练；
对比对齐机制的实现与消融分析；
因果驱动条件流匹配的端到端训练；
因果效应提取与生物学验证；
综合实验评估与基线方法全面比较 &
完成第二个研究点的算法设计与实验验证；
形成第二篇可投稿的研究论文 \\
\hline
2027年4月--2027年5月 &
学位论文撰写与修改；
论文投稿与答辩准备；
开源工具包的完善与文档撰写 &
完成硕士学位论文并通过答辩；
完成1-2篇学术论文投稿；
正式发布开源工具包 \\
\hline
\end{tabular}
\end{table}

\nocite{*}% 使文献列表显示所有参考文献（包括未引用文献）
%---------------------------------------------------------------------------%
